\documentclass[tikz, border=10pt]{standalone}
\usepackage{pgfplots}
\pgfplotsset{compat=1.18}

\begin{document}
\definecolor{maincolor1}{RGB}{1,82,142}
\definecolor{maincolor2}{RGB}{21,182,184}
\definecolor{mygrey}{rgb}{0.41, 0.41, 0.41}
\begin{tikzpicture}
\begin{axis}[
title={$f(x) = x^2 - 3.2x + 3 \quad (n = 6)$},
xlabel={$x$},
ylabel={$f(x)$},
xmin=-0.2, xmax=3.2,
ymin=-0.1, ymax=3.3,
axis lines=left,
grid style={line width=.1pt, draw=gray!20},
major grid style={line width=.2pt, draw=gray!50},
minor tick num=1,
legend pos=north west,
legend style={font=\small, fill=white, fill opacity=0.9, draw opacity=1},
width=12cm,
height=8cm
]

% 1. Región sombreada del intervalo de incertidumbre [1.0, 2.0]
\addplot[draw=none, fill=mygrey, opacity=0.1] coordinates {
(1.0, -0.1) (1.0, 3.3) (2.0, 3.3) (2.0, -0.1)
};

% 2. Curva continua f(x)
\addplot [
domain=0:3, 
samples=100, 
color=maincolor1, 
thick
] {x^2 - 3.2*x + 3};

% 3. Líneas verticales proyectadas a la malla
\pgfplotsinvokeforeach{0, 0.5, 1.0, 1.5, 2.0, 2.5, 3.0}{
\draw[gray!50, dashed] (axis cs:#1, -0.1) -- (axis cs:#1, {#1^2 - 3.2*#1 + 3});
}

% 4. Puntos evaluados en la malla de búsqueda
\addplot [
only marks, 
mark=*, 
mark size=3pt, 
color=maincolor2
] coordinates {
(0.0, 3.0)
(0.5, 1.65)
(1.0, 0.80)
(1.5, 0.45)
(2.0, 0.60)
(2.5, 1.25)
(3.0, 2.40)
};

% 5. Mínimo analítico exacto (1.6, 0.44)
\addplot [
only marks, 
mark=diamond*, 
mark size=4pt, 
color=maincolor2,
thick
] coordinates {(1.6, 0.44)};

\end{axis}
\end{tikzpicture}
\end{document}